% Порождено tools/gen_map.py, не править вручную.
\chapter{Соответствие глав и уроков Logra}
\label{app:logramap}

В скобках указана книга: «Схемы» и «Verilog» --- книги вкладки \emph{Обучение}, «Курс» --- её управляемые уроки. Таблицы показывают, какой урок
и какой пример (вкладки \emph{Код} и \emph{Конструктор}) соответствуют главам книги. Названия сверены с исходным
кодом Logra при сборке этого приложения.

\section*{Главы по порядку книги}
\begin{longtable}{@{}p{0.24\textwidth}p{0.38\textwidth}p{0.26\textwidth}@{}}
\toprule Глава & Урок Logra (книга) & Пример (\code{id}) \\ \midrule \endfirsthead
\toprule Глава & Урок Logra (книга) & Пример (\code{id}) \\ \midrule \endhead
Что такое ПЛИС & Что у нас в руках (Курс); 1. Осмотр мастерской (Курс) & \code{blink} \\
Среда Logra & Ты в мастерской цифровых схем (Курс); Куда ты попал (Курс) & \code{blink} \\
Сигналы & Бит: ноль и единица (Курс); Логический элемент (Курс); Число из четырёх бит (Курс) & \code{button}, \code{not}, \code{blink} \\
Системы счисления и коды & Двоично-десятичный код (Схемы); Код Грея блоком (Схемы); Бит чётности (Verilog) & \code{gray}, \code{parity}, \code{binary-bcd}, \code{hamming-encode} \\
Логические элементы и булева алгебра & И-НЕ из блоков (Схемы); Закон де Моргана (Схемы); XOR из четырёх И-НЕ (Схемы) & \code{gates}, \code{inverse-gates}, \code{demorgan} \\
Основы Verilog & Первый провод (Курс); Склей биты (Курс); Выбор по условию (Курс); Первая ячейка памяти (Курс); Модуль внутри модуля (Курс) & \code{gates}, \code{not} \\
Комбинационные схемы & Два из трёх (Схемы); Мультиплексор из НЕ, И, ИЛИ (Схемы); Голосование «два из трёх» (Verilog) & \code{logic}, \code{majority} \\
Минимизация логических функций & Равенство: XNOR (Схемы); ИЛИ на три входа (Схемы) & \code{implication}, \code{consensus} \\
Дешифраторы и шифраторы & Дешифратор (Схемы); Приоритетный шифратор (Схемы); Дешифратор 3 в 8 (Verilog) & \code{decoder}, \code{priority} \\
Мультиплексоры и демультиплексоры & Мультиплексор 2→1 (Схемы); Мультиплексор 4→1 (Схемы); Демультиплексор (Схемы) & \code{mux}, \code{mux-four}, \code{demux-four} \\
Арифметические узлы & Полусумматор из вентилей (Схемы); Полный сумматор (Схемы); Вычитание (Схемы); Меньше: A < B (Схемы); Четырёхразрядный сумматор (Verilog) & \code{half-adder}, \code{full-adder}, \code{comparator} \\
Триггеры & Память на один бит (Курс); SR-триггер: установить и сбросить (Схемы); D-триггер: запомнить бит (Схемы); Детектор фронта: один импульс на нажатие (Схемы) & \code{edge-toggle} \\
Регистры и счётчики & Регистр на четыре бита (Схемы); Сдвиговый регистр блоком (Схемы); Счётчик одним блоком (Схемы); Счётчик с паузой (Схемы) & \code{shift}, \code{counter}, \code{johnson} \\
Синхронные схемы, автоматы и память & Синхронизатор: привязать кнопку к такту (Схемы); Делитель частоты на 2 (Схемы); ПЗУ из массива (Verilog) & \code{async-toggle-sync}, \code{button-debounce}, \code{divide-by-three} \\
Проектирование автоматов и очереди FIFO & Собери светофор (Курс); Детектор 101 (Схемы); Кодовый замок (Схемы); Выход Мура (Схемы) & \code{traffic}, \code{vending-fsm}, \code{sequence-1011}, \code{fifo-four} \\
Что внутри ПЛИС & Сколько длится один период? (Курс) & \code{logic} \\
Плата Sipeed Tang Nano 20K & Кнопка на настоящей плате (Курс); Светодиод горит нулём (Verilog) & \code{button} \\
От текста до работающей схемы & Привет! (Курс); Светодиод и кнопка вместе (Verilog) & \code{button} \\
Электрические основы: КМОП, выходы и шины & Светодиод и ограничение тока (Схемы); Тумблер (Схемы) & \code{button} \\
Цифра и аналоговый мир: ЦАП, АЦП, ШИМ & ШИМ: яркость числом (Verilog) & \code{pwm}, \code{breathing-pwm}, \code{sigma-delta} \\
От схемы к Verilog: узлы теоретических глав & Полный сумматор: учитываем перенос (Verilog); Счётчик с загрузкой (Verilog); Мур и Мили: когда выход реагирует (Verilog) & \code{full-adder}, \code{decoder}, \code{traffic}, \code{fifo-four}, \code{hamming-syndrome} \\
Практика: семь проектов & Бегущий огонёк (Verilog); Мигание: делим частоту на 8 (Verilog); Нажал: переключил (Verilog) & \code{gates}, \code{blink}, \code{chaser}, \code{button-debounce}, \code{traffic}, \code{breathing-pwm}, \code{mux} \\
Симуляция & Сколько длится один период? (Курс) & \code{blink} \\
Тактирование и временной анализ & Синхронизатор из двух триггеров (Verilog); Сброс сразу или по такту (Verilog) & \code{async-toggle-sync}, \code{button-debounce} \\
UART, SPI, I2C & --- & \code{uart-tx}, \code{uart-rx}, \code{spi-loopback}, \code{i2c-pattern} \\
Память SDRAM & ОЗУ 8×4: запись по такту (Verilog); ПЗУ из массива (Verilog) & \code{ram-write-read}, \code{small-rom} \\
Собираем процессор & Мини-АЛУ на четыре бита (Verilog) & \code{alu}, \code{tiny-cpu} \\
Задания для самостоятельной работы & Секундомер по кнопке (Verilog); Охрана: два датчика из трёх (Verilog) & \code{stopwatch} \\
\bottomrule
\end{longtable}

\section*{Сквозной проект «Охранная панель»}
\begin{longtable}{@{}>{\raggedright\arraybackslash}p{0.20\textwidth}>{\raggedright\arraybackslash}p{0.24\textwidth}>{\raggedright\arraybackslash}p{0.26\textwidth}>{\raggedright\arraybackslash}p{0.17\textwidth}@{}}
\toprule Этап & Что собираем & Урок или пример & Критерий приёмки \\ \midrule \endfirsthead
\toprule Этап & Что собираем & Урок или пример & Критерий приёмки \\ \midrule \endhead
Сигналы & Датчик двери как бит; светодиод «тревога» на кнопке & Бит: ноль и единица (Курс), \code{button} & LED светится при нажатой кнопке \\
Логические элементы & alarm = дверь И (ночь ИЛИ охрана) & И на три входа (Схемы), \code{gates} & таблица на Диаграммах совпала с формулой \\
Комбинационные схемы и минимизация & Три датчика, голосование «2 из 3», минимальная форма & Голосование «2 из 3» с охраной (Схемы); Охрана: два датчика из трёх (Verilog), \code{majority} & все 8 комбинаций; число LUT не больше одной \\
Дешифраторы и мультиплексоры & Выбор зоны и показ её номера на семи сегментах & Цифра на семи сегментах (Схемы), \code{decoder}, \code{mux} & номер зоны виден на индикаторе \\
Арифметические узлы & Порог: уровень шума A > B & Меньше: A < B (Схемы), \code{comparator} & тревога только выше порога \\
Коды & Код замка в BCD и бит чётности при передаче & Проверка чётности слова (Схемы), \code{parity} & ошибка в одном бите обнаруживается \\
Триггеры & Запоминание тревоги, сброс ключом & Сигнализация с запоминанием (Схемы); Сигнализация с памятью (Verilog), \code{edge-toggle} & тревога держится после ухода сигнала \\
Регистры и счётчики & Счётчик времени тревоги, делитель для писка & Удлинитель импульса (Схемы), \code{counter} & сирена выключается через заданное число тактов \\
Автоматы & ОХРАНА\_ВЫКЛ, ОХРАНА, ТРЕВОГА; ввод кода 1-0-1 & Кодовый замок (Схемы); Кодовый замок: 1, 0, 1 (Verilog), \code{door-controller} & неверный код не снимает охрану \\
Электрика и ШИМ & Зуммер через ШИМ, подтяжка кнопок, плата & Зуммер (Схемы); ШИМ: яркость числом (Verilog), \code{pwm} & работает на плате Tang Nano 20K \\
\bottomrule
\end{longtable}
